\begin{otherlanguage}{english}
\noindent\textbf{Coefficients, monodromy, and the search for a usable fibration.}

These pages of Chapter~I begin with a correction to an overly strong use of cofinality and end with a programme for a theory of fibrations in the category of small categories. Their continuity lies in a single question: which coefficients must a comparison preserve if that comparison is to remain useful after another operation? Preservation of constant coefficients is enough to define a class of weak equivalences. It is not, in general, enough to carry out base change or to compose fibrations. The passage from constants to locally constant coherent coefficients is therefore the principal mathematical movement of the text.

These are pages~34--74 of the edition, manuscript pages~60--126. The text follows the editors' electronic edition \citep{Grothendieck1990Derivateurs}. Manuscript pages~60--66 also close the first part of the chapter, where their commentary, claims and readings are given. Printed and manuscript numbers are kept distinct throughout. No assertion in the commentary silently corrects a mathematical mistake in the French source: genuine transcription errors are corrected, whereas errors already present in the edition remain visible and are explained where they affect the argument.

We use the manuscript's convention that an object of $\D(X)$ has underlying diagram $X^{\mathrm{op}}\to\D(e)$. Thus the value of a right Kan extension $f_*F$ at $y$ is calculated on $X/y$, and the value of the left Kan extension $f_!F$ is calculated on $y\backslash X$. The symbols $H^*_{\D}$ and $H_*^{\D}$ denote objects in the coefficient category, not merely lists of cohomology or homology groups. This convention is responsible for several apparently reversed statements when one compares the manuscript with modern covariant-diagram treatments; the variance distinction is explained explicitly in \citet[Warning~2.2]{GrothPontoShulman2014}.

\medskip\noindent\textbf{From constant equivalences to base change.}

For $f:X\to Y$, with terminal projections $p$ and $q$, a $\D$-equivalence means that restriction induces bijections
\[
 \operatorname{Hom}_{\D(Y)}(q^*M,q^*N)
 \longrightarrow
 \operatorname{Hom}_{\D(X)}(p^*M,p^*N)
\]
for all $M,N\in\D(e)$. Adjunction and Yoneda identify this with invariance of constant-coefficient homology or cohomology whenever the relevant Kan extensions exist. It does not say that $f^*$ is fully faithful on all of $\D(Y)$. The distinction is already decisive for an endpoint inclusion in the walking arrow.

Theorem~1 on manuscript page~67 gives the stronger cofinality test actually needed: cohomological invariance for every coefficient on $Y$ is equivalent to asphericity of the comma categories $y\backslash X$. The proof uses the counit $f_!p^*M\to q^*M$. One adjunction reads it as a cohomology comparison; pointwise evaluation reads it as constant homology of a comma category. This explains the apparent use of homology inside a cohomological theorem without adding a missing hypothesis to a valid argument.

The proper and smooth base-change results then have a precise local content. Properness requires the categories of lifts of every arrow $f(x)\to y$ starting at $x$ to be $\D$-aspherical. The requirement involves both objects and their morphisms, rather than mere nonemptiness of the set of lifts. A Grothendieck opfibration satisfies it because an opcartesian lift is an initial object of the corresponding lift category. The dual account gives smoothness. The commentary supplies the omitted implications among the endpoint, lifting, and base-change tests in Theorem~2, while preserving the edition's mistyped objects and cross-references as source readings.

This vocabulary retains the geometry of proper and smooth base change, but defines a property of functors relative to a coefficient theory. The later theory of asphericity structures provides an appropriate comparison, rather than an identification with topological proper maps \citep[Theorems~1.2, 4.13, and~5.4]{Maltsiniotis2005Asphericite}. The intervening construction $\Phi(X)$, now called the twisted-arrow category, gives natural comparisons between $X$ and $X^{\mathrm{op}}$. Both endpoint projections are opfibrations with contractible fibres. Their universal equivalence proves that taking opposites preserves the equivalence class detected by $\D$. This is a comparison of homotopy information, not an equivalence of the original categories.

\medskip\noindent\textbf{What the counterexamples distinguish.}

From manuscript page~76 onward, the word ``fibration'' changes meaning as the argument develops. Its first version asks that the actual fibre and its comma-category model have the same constant-coefficient theory and that transport between comma categories preserve that theory. Requiring this after every base change is an additional condition. We exhibit both a two-object example and an example involving only finite posets where the first condition holds but fails after a strict base change.

The more substantial failure concerns the proposed implication that a universal fibration must preserve every $\D$-equivalence of bases. For rational complexes, a base can be cohomologically indistinguishable from a point while still carrying nontrivial local systems. Grothendieck's reflected-circle bundle over the circle makes the point concretely. The degree-one cohomology of its fibres is a sign local system on the base circle; its derived cohomology has differential multiplication by $-2$, hence is zero over $\mathbb Q$. The projection from the Klein bottle is consequently a rational constant-coefficient equivalence although its circle fibres are not rationally aspherical. In a fibre square the tensor square of the orientation local system becomes untwisted, and the extra cohomology reappears. The commentary calculates the full complexes and the comparison maps, including the degree-zero term omitted from the source's abbreviated group-cohomology calculation.

The finite-group version replaces this computation by the averaging idempotent. It also shows why the proposed extension to truncated homotopy theories needs care: already the point mapping into $BC_2$ is an equivalence for the represented set-valued derivator, whereas locally constant $C_2$-sets retain the action. The corresponding assertions about complete equivalences cannot be inferred from connectedness.

One further correction is necessary in the proposed finite-poset model. The action on $S^2$ is antipodal and the action on $S^1$ is a reflection. Their product acts trivially on the top class of degree three. Hence
\[
 H^*((S^2\times S^1)/C_2;\mathbb Q)
 \cong\mathbb Q\quad\text{in degrees }0,3,
\]
and vanishes in other degrees. The base $\mathbb{RP}^2$ is rationally aspherical, but this total space is not. The example still disproves invariance under inclusion of a fibre; it does not supply the additional aspherical-total-space claim. An exact rational chain calculation for the thirteen-object base and fifty-two-object total poset independently confirms these Betti numbers. The orbit-category relations printed before this construction also require correction, which is explained separately in the commentary.

\medskip\noindent\textbf{The final distinctions and their present interpretation.}

The later definitions separate constant-coefficient equivalence, invariance of cohomology with every locally constant coefficient, and full faithfulness or equivalence on the entire category of locally constant objects. The last two are called $1$- and $2$-complete equivalence. A closed internal-Hom formalism identifies cohomological equivalence with $1$-complete equivalence, under the stated compatibility assumptions. It does not identify ordinary constant-coefficient equivalence with either of them. The missing word ``cohomological'' in the sentence on manuscript page~122 is therefore a material slip.

The nerve-invariance assertions eventually formulated as the $b'$ variants of \textbf{Der~3} have a precise later answer. \citet[Theorem~1.3 and Corollary~3.8]{Cisinski2009LocallyConstant} proves invariance of locally constant coefficients under weak equivalences of indexing nerves, first in the model-category setting and then for derivators. In the space-valued case these coefficients can be read as coherent functors out of the fundamental infinity-groupoid. A weak equivalence changes neither that groupoid nor its coefficient categories. This establishes the strong form of the lemma about $\HOT$ that Grothendieck had deferred; it does not collapse the distinctions for a fixed linear coefficient theory. The minimal-localizer theorem places the ordinary weak-equivalence class inside the broader framework that motivated these pages \citep[Theorem~2.2.11]{Cisinski2004Minimal}.

Two developments in the commentary on these pages address unfinished questions directly. First, for a universally weak cohomological fibration, the formula for $f_*f^*F$ holds for every coefficient $F$ on the base, with no local-constancy assumption. A factorization category for each arrow $z\to y$ reduces the left mate of the desired comparison to constant homology. The dual proof applies to $f_!f^*F$. As a consequence, a universal $\D$-equivalence induces fully faithful restriction on the whole coefficient category. The required closure of the diagram domain is stated explicitly. These arguments use the factorization method underlying homotopy exactness \citep[Theorem~3.8]{GrothPontoShulman2014}, and the full-faithfulness conclusion for universal weak equivalences of nerves also follows from the modern homotopical epimorphism criterion \citep[Theorem~6.10]{GrothStovicek2016Trees}.

Second, the intermediate $\D$-fibrations are not closed under composition. Over $\mathbb Q$, take the opfibration over $[1]$ with transition $e\to BC_2$ and then the covering determined by the regular $C_2$-set. Both factors satisfy the universal constant-coefficient condition. Their composite has transition from two discrete points to a contractible action groupoid, and fails already in degree-zero cohomology. The stronger cohomological and homological fibrations on manuscript page~121 control locally constant coefficients on the total category and are closed under composition. A separate example with transition $BC_2\to e$ shows that this stronger relative theory has examples beyond $\HOT$-fibrations. The proposed $2$-complete variant also fails to compose for the represented set-valued derivator: a $C_2$-valued cocycle on the tetrahedral boundary face poset gives two functors between finite categories whose composite fails the required test. The obstructed restriction includes sets as trivial $C_2$-sets; its failure of essential surjectivity distinguishes this counterexample from the $1$-complete case, which composes under the internal-Hom hypotheses.

The mathematical status is thus stated at the level of the actual assertions. The main cofinality and base-change theorems hold with their conventions and domain hypotheses made explicit; the false implications have concrete counterexamples; several stated questions admit the supplied proofs or the cited later theorems. A failure to locate a particular construction in the sources checked is recorded as such, never converted into a claim that it is an open problem or a new result. The label \ours\ marks proofs and constructions given in the commentary; it carries no priority claim.
\end{otherlanguage}
